\section{The Western Tradition Reviewed}
\begin{quotex}
In which we review some prior posts on the western tradition and demonstrate the inner continuity from the Vedics through the Greeks to the Medievals. We also articulate the essence of idealism and the nature of the priest.
\end{quotex}

We have several times pointed to the three great Indo-European civilizations: the Hindu, the Greco-Roman, and the Medieval. That is its depiction from the outside, but we have alluded to their interior connection, which is more important. Specifically, that means that they appear as three forms of the same inner spirit. Interestingly enough, \textbf{Vladimir Solovyov} approaches them from their interiority, identifying the same three forms as did Evola. The inner relationship of these forms will be taken up more fully at a later time.

There may sometimes be the feeling that the ``transcendental unity of religions'' represents some sort of lowest common denominator. Solovyov indicates the opposite when he writes:

\begin{quotex}
From the religious point of view, the goal is not a minimum but a maximum of positive content. The richer, the more alive and concrete a religious form is, the higher it is. The perfect religion is not one that is equally contained in all religions (the indifferent foundation of religion); the perfect religion is one that possesses and contains within itself all religions (the complete religious synthesis). The perfect religion must be free from all limitation and exclusiveness, not because it is deprived of every positive particularity and individuality, but because it contains in itself all particularities and is not exclusively attached to any one of them, possesses all of them, and is therefore free from all of them. The dark fanaticism that holds to a single particular revelation a single positive form, denying all others, and the abstract rationalism that disperses the whole essences of religion into a fog of indeterminate concepts, fusing all religious forms into a single empty, impotent, colorless generality, are equally repugnant to the true conception of religions.
\end{quotex}


Solovyov points out the Christian metaphysics is compatible with that tradition, and that is why they could speak of ``Christians before Christ''. He even points that the Egyptian theosophy of \textbf{Hermes Trismegistus} is a part thereof. He explains it this way:

\begin{quotex}
The originality of Christianity lies not in its general views but in positive facts, not in the speculative content of its idea but in its personal incarnation.
\end{quotex}


But those positive facts are outside the scope of this discussion, so we are content to point out the identity of its speculative, or metaphysical, content with what preceded it.

\paragraph{Roman High Priest}


For over 2700 years, the Roman High Priest has been the spiritual center of the Ecumene. Recently, a new one was selected. Although I am not too concerned with the specifics of the who, what, and where, it is nevertheless important that the seat be filled. Nevertheless, the reactions of the secular media was interesting.

First of all, they were all impressed that Pope Francis is ``for the poor'', something they derive from his non-ostentatious lifestyle. That is something they can grasp, since it is one of the corporeal works of mercy. These works are acceptable to them because they are based on the natural principle of ``do unto others what you would have them do unto you.'' Obviously, a hungry man, woman, or child should be fed, when possible. However, that does not mean the creation of a permanent underclass that does not feed itself when possible. That is the Marxist principle of ``to each according to his needs'', which is not the same thing.

Of course, man does not live by bread alone, so the spiritual works of mercy are more important. Here the secular media falls away. The first of them is ``to instruct the ignorant''; the sophists of the fourth estate do not consider themselves ignorant, and so they resist such instruction. They seem to tolerate it in Francis, since they presume it is merely a relic from bygone times and will eventually be forgotten in the general progress of humanity.

Ultimately, none of that matters, since the true idea of priesthood itself is what has been forgotten, not only among the secularists, but it seems even among the priests. \textbf{Guido De Giorgio} seems to capture that essence:

\begin{quotex}
Priests are not men, they are vowed, they are consecrated, they have nothing apart from God, since they chose God and cannot, and must not, forget it, but not even those who get close to them can forget it. The priest has no family, no ties, he is excluded from the world, and excluded voluntarily, a denier of ``ordinary life'' for having chosen the ``extraordinary life'', that adheres totally to the divine truths and acts on men negatively, in order to distract them from the fallacies of the world and to make them understand what they lose while winning the game of life. If he is connected to few living in the world outside the world, in continuous contact to sacred truths, in a way that they can, without apparent sacrifice, pass unencumbered, that cannot and must not constitute a good formula for the majority who, if it admits the world, subtracts God.
\end{quotex}


\paragraph{Magical Idealism}


Yet, it is insufficient to dwell on the speculative content alone with actualizing its contents. Metaphysics is a science, not a system of rational discourse. In \emph{Essays on Magical Idealism}, Evola points to some of the roots of his thought, among which is Giovanni Gentile, who was the foremost Italian philosopher at that time. Evola points out that Gentile limits himself to the disputes of university professors, when his system of thought was quite close to a method of self-realization. Evola writes:

\begin{quotex}
Modern idealism can be defined this way: a deep need toward an absolute self-realization, which however the I does not acquire immediately in its interiority, but it simply \emph{knows}, comprehending it from the outside in the phenomenon that it produces in the abstractly rational order. In Giovanni Gentile this situation appears in a particularly clear way: in him the effort of encompassing and dominating the whole of the world in an immanent principle reaches its perfection; but, on the other hand, this principle remains a simple ideal entity, it is the ``transcendental I'' that we criticized and yet it expresses only the pallid reflection of that deep individual power that existed for example in Michelstaedter. If Gentile could really call ``I'' the ``pure act'' of his rationalism, then one would find him appearing not as the university professor, whose ``actuality'' has as its goal the reform of the scholastic program, but as that cosmic centrality that the esoteric teaching points to, as examples, the type of the \emph{rishi}, the \emph{yogi}, the \textbf{Christ}, and the \textbf{Buddha}.
\end{quotex}



\flrightit{Posted on 2013-03-25 by Cologero}

\begin{center}* * *\end{center}

\begin{footnotesize}
\begin{sffamily}
\texttt{Synodius on 2013-03-26 at 12:25 said:}

What is disturbing about this new pope is that under the mask of humbleness he is giving up papal symbols and traditions which represented the spiritual authority. He finds it hard to perform even such a simple act as papal blessing in order to not offend anyone. It is not Francis personally that matters, but the fact that the revolution is fully in charge since the last council. It seems that this pope may give the Church the final blow with irreversible democratization, wild ecumenism and support of socialism

\hfill

\texttt{George Goerlich on 2013-04-01 at 22:11 said:}

I noticed the same as Synodius. It appears he is symbolically assuming the role of the ``anti-Pope''. Not in the sense of Satan, but in the sense of not-being a Pope. A Pope who denies the title, symbols, and position of a Pope.

\end{sffamily}
\end{footnotesize}

